% 正赛文稿：反方。环节标题按 recruit-3v3.md「各持方文稿标题」原样使用。
\documentclass[script,con]{debate}
\motion{“磕CP”热潮折射出人们对爱情的更深信仰 / 更深怀疑}
\stance{反方}{更深怀疑}
\meta{赛制}{招新赛 3v3}
\meta{口径来源}{备赛文档 第十节 反方口径表}
\begin{document}
\debatetitle
\begin{thesisbox}[全场一句话立论]人们只敢把爱放在屏幕那头、结局写好的地方；信的是别人的爱，疑的是自己的爱靠不靠得住，疑在底下。\end{thesisbox}

\begin{kit}{1. 反一开篇立论}{套件 · 组装后 822–835 字 / 180 秒}
\assembly{开头（任选 1） → 立论主体}
\begin{pick}{开头}{1}
\begin{module}{如果质询拿到了正一承认“嗑CP的人不进入那段关系”}
谢谢主席。刚才质询里，对方一辩承认了一件事：嗑CP的人，不进入那段关系。请评委记住这句话，我方的立论就从这里开始。
\end{module}
\begin{module}{如果质询里正一回避了“既然信，为什么不去谈”}
谢谢主席。刚才质询里，我方问对方一辩，既然信爱情，为什么不去谈。这个问题，对方没有正面回答。我方的立论就从这个问题开始。
\end{module}
\begin{fallback}{质询没有拿到明确成果}
谢谢主席，各位好。同一股热潮，大家看到的是甜。我方想请评委，看一看这份甜是放在什么地方的。
\end{fallback}
\end{pick}
\begin{fixed}{立论主体}
先把几个词说清。嗑CP，是以旁观者的身份，给虚构角色或者真实人物配对，想象并期待他们之间的浪漫关系，从中得到情感满足。对爱情的信仰，是相信真正的爱情存在，值得期待和投入；怀疑，是不相信它存在，或者不相信它靠得住、值得投入。

热潮里两种心态都有，这一点我方承认。所以“更深”，比的是谁在底下。判断标准是：哪一种心态更能解释热潮的核心特征，也就是人们嗑的是什么样的爱，为什么用旁观的方式去爱，并且它是另一种心态的来源。在这个标准下，对方要证明，嗑CP的动力来自相信真爱存在；我方要证明，嗑CP的动力来自不信真爱靠得住，热潮里的甜，只在安全距离内成立。

我方第一点，信在远处。

嗑CP，从形式上就定了一件事：那段爱，你进不去，也不用进去。你在屏幕这头，爱在屏幕那头。热潮越大，越多的人把爱放在一个自己不必进场的位置。那他们自己的恋爱呢？共青团中央的课题组二〇二一年调查了两千九百多名十八到二十六岁的未婚城市青年，问将来会不会谈恋爱，一成多说不谈，还有超过四分之一说不确定。

人民大学的许冠文研究嗑CP，他在采访里说，CP粉里流传着一句话：“虽然我不谈恋爱，但是我要看我的CP谈恋爱。”他自己也好奇，为什么她们偏向消费爱情的产品，而不是自己去实践爱情。请评委听清那句话的结构。爱情会发生，但发生在别人身上；落到自己，是“虽然我不”。一份只对别人成立的信，底下就是疑。

我方第二点，要的是保险。

真实的爱情一定带着风险：会变，会散，会受伤。嗑CP恰好把这些全拿掉了。许冠文的论文里，一位受访者说：如果在现实生活中去体验这种爱情，自己很难不发疯。许冠文的解释是，这种爱情带来的风险是非常大的。记者问他，这样是不是更安全，他说：对，你不用担心他们会解散，因为你已经给他们画好了终点。

说白了，结局提前写好，才敢投入。这背后是一个判断：真实的爱，结局不会自己好。对爱靠不靠得住，嗑的人已经投了票。

所以我方全场都会请对方回答一个问题：如果人们真信爱情，为什么要把它放在一段自己碰不到、结局早就写好的关系里？

我方不是在指责嗑CP的人，只是想听清他们没说出口的那半句。热潮越甜，越说明\stress{现实里有人不敢尝}。谢谢。
\end{fixed}
\end{kit}
\begin{contingency}
\ifthen{正方把“信仰”定义成“只要喜欢爱情故事就算信”}{那看言情小说的人全都信爱情，这个定义让题目失去意义；请回到中立定义：信仰包括相信它靠得住、值得投入。}
\ifthen{正方说“嗑的是专一和双向，所以是信”}{想象里的标准不用兑现，所以才敢拉满；判准要问的是为什么用旁观的方式去爱。}
\ifthen{正方拿出“将近八成想谈恋爱”}{想谈不等于相信；想而不去，正是我方说的那份疑。}
\ifthen{论点一被打成“旁观不等于怀疑”}{收缩到论点二：画好终点才敢投入，这是对靠不靠得住的判断。}
\end{contingency}

\begin{stage}{2. 反一被质询防守预案}{正二质询，90 秒双边计时}
\textbf{原则}：先答是或否，再守一句前提；对方问内容，我方拉回形式与来源。
\begin{qa}
\qarow{大家嗑得最多的，是专一的一对还是互相算计的一对？}{专一的。想象里的专一不用兑现，所以才敢要；我方看的是为什么只敢在想象里要。}
\qarow{不信爱情的人，会为专一的一对上头吗？}{会。因为那一对不会让他受伤。上头的是安全的专一。}
\qarow{没脱单排第一的是圈子窄，对吗？}{那是全体单身青年的调查；圈子窄的人也可以去认识人，嗑CP的人选了屏幕。}
\qarow{看球的人不上场，能说他不信足球吗？}{看球的人没说过“球场太危险，我只看”。嗑的人说了：现实里体验这种爱，很难不发疯。}
\qarow{那句“虽然我不谈恋爱”，是受访者说的吗？}{是许冠文转述的 CP 粉圈里流传的一句话，他在采访里原话引用。}
\qarow{您方承认热潮里也有信吗？}{承认。信在屏幕那头；我方比的是谁在底下。}
\end{qa}
\end{stage}

\begin{stage}{3. 反二质询正一问题链}{90 秒，双边计时}
\textbf{总目标}：全场第一次质询，为反一立论开头取材：对方承认嗑CP的人不进入那段关系；对方回避“既然信，为什么不去谈”。
\begin{chain}{链一：形式链}{拿到“不进入那段关系”的承认}
\q{嗑CP的人，自己进入那段关系吗？}{答不进入：下一问；答进入：怎么进入？}
\q{他们选的，是一段自己进不去的爱，对吗？}{答对：收束；答不对：请举一对嗑的人能进去的CP}
\close{感谢，对方承认嗑CP的人不进入那段关系。}
\end{chain}
\begin{chain}{链二：行动链}{逼出“为什么不去谈”}
\q{您方说信爱情，那将近八成想谈的人，谈了吗？}{答没有：下一问；答有的谈了：那没谈的呢？}
\q{既然信，您方说得出不去谈的理由吗？}{答圈子窄：圈子窄就只能嗑CP吗？答不回答：复述一次，换下一问}
\q{不敢去，算不算不信它靠得住？}{答不算：请评委记住对方把不敢和不信分开；答算：收束}
\cut{好，我换个问法：想要而不敢，您方叫它什么？}
\close{感谢，“既然信为什么不去谈”，请对方后面继续回答。}
\end{chain}
\end{stage}

\begin{stage}{4. 反三对辩要点}{双方各 90 秒}
\textbf{任务}：接住正三开的“内容与形式”战场，把它拉到形式与来源。每次先回一句，再问一句。
\begin{points}{进攻点（4–5 个，每个 15–20 秒）}
\item 专一、双向都放在一段自己进不去的关系里。信，为什么只敢信在那里？
\item 想象里的标准不用兑现，所以才敢拉满。标准高，恰恰是因为不用付钱。
\item 嗑的人自己说：现实里体验这种爱，很难不发疯。这是信，还是怕？
\item 画好终点才敢投入，说明不信结局会自己好。
\item 更深比的是来源：是因为不敢，才去嗑；还是因为想要，才不去？
\end{points}
\begin{points}{备用点（6–8 个）}
\item 对方说旁观不等于怀疑：看球的人没说球场太危险；嗑的人说了。
\item 对方说圈子窄：圈子窄的人也能去认识人，他们选了屏幕。
\item 对方说“双向奔赴”进了流行语：流行的是一个词，不是一次行动。
\item 对方说将近八成想谈：想和信之间，差的就是敢不敢。
\item 对方说怕说明当真：信它是真的，却不信它靠得住，这正是中立定义里的怀疑。
\item 对方说代餐说明饿：饿却只吃代餐，说明不信正餐吃得到。
\item 对方说宁缺毋滥：落到现实，就是缺。
\end{points}
\begin{points}{对方最可能问的三个问题 → 回应}
\item 不信爱情的人，为什么嗑专一？ → 专一放在一段不会伤到自己的关系里，才安全。
\item 旁观就是怀疑？ → 不是所有旁观，是说出了“现实太险”的旁观。
\item 你方数据问的是嗑CP的人吗？ → 共青团那份是全体未婚青年，说明背景；许冠文的研究对象就是CP粉。
\end{points}
\end{stage}

\begin{kit}{5. 反二申论}{套件 · 组装后 388–403 字 / 90 秒}
\assembly{开头 → 拆正方立论（任选 1） → 推进论点一 → 收尾}
\begin{fixed}{开头}
谢谢主席，各位。我这一篇做两件事：先把正方的立论放到判准下看，再把我方第一点往前推一步。
\end{fixed}
\begin{pick}{拆正方立论}{1}
\begin{module}{如果正方立论的主干是“嗑的是理想，所以是信”}
正方的立论说，大家嗑的是理想的爱，所以是信。这里缺了一步。想象里的标准，不用兑现，所以才敢拉满。叶公对龙的标准也很高，真龙来了，他跑了。正方要证明的不是人们想要什么样的爱，而是人们相信这样的爱会来。标准再高，放在一段自己进不去的关系里，它说明的是向往，不是相信。向往和相信之间，隔着的就是一个“敢”字。
\end{module}
\begin{module}{如果正方立论的主干是“疑的只是遇见，不是爱情”}
正方的立论说，年轻人疑的只是遇不遇得到，不是爱情本身。这一步分得太干净了。请看中立定义：怀疑包括不相信它靠得住。一份爱，如果扛不住圈子窄，扛不住热搜上的负面新闻，那它在人们心里就是靠不住的。想谈，是一回事；相信谈了会好，是另一回事。就算拿出想谈的数据，也只证明了前一件；后一件，一个字都证明不了。想要却不敢，这正是怀疑最常见的样子。
\end{module}
\begin{fallback}{正方立论的主干不在以上两条时}
正方的立论，请评委放回判准里看。判准问的是两件事：人们嗑的是什么样的爱，为什么用旁观的方式去爱。一份立论如果只讲嗑的内容，不讲为什么偏偏用旁观的方式，为什么偏偏把爱放在屏幕那头，它就只回答了一半。而没回答的那一半，恰恰决定了底下是什么。一个人想要什么，和他敢不敢要，是两回事；题目问的“更深”，就藏在敢不敢里。
\end{fallback}
\end{pick}
\begin{fixed}{推进论点一}
我方第一点，信在远处，今天再补一层。两位传播学者 Horton 和 Wohl 在一九五六年就描述过这种隔着屏幕的关系：观众几乎不用承担义务，也不用付出努力，随时可以抽身；这种关系是单向的，不会一起往前走。你想想，这正是嗑CP最舒服的地方。可一个人把爱放在一个不需要付出、随时能走的位置，说明他还不敢把爱放在需要付出、不能随便走的位置上。许冠文自己也问过同一个问题：为什么她们偏向消费爱情的产品，而不是自己去实践爱情？
\end{fixed}
\begin{fixed}{收尾}
说到底，请正方回答一个问题：如果人们真信爱情，为什么只敢把它放在随时可以抽身的地方？
\end{fixed}
\end{kit}

\begin{stage}{6. 反二被质询防守预案}{正三质询，90 秒双边计时}
\textbf{原则}：申论里的 Horton 和 Wohl 下一分钟就会被问，出处和边界都要背。
\begin{qa}
\qarow{您方说怀疑包括不信它靠得住，对吗？}{对，这是双方共用的中立定义。}
\qarow{怕受伤的人，是把爱当真，还是当假？}{当真。信它是真的，却不信它靠得住，就是定义里的怀疑。}
\qarow{Horton 和 Wohl 说对多数人是补充，对吗？}{对，他们写的是一九五六年的电视观众。补充还是替代，要看人们自己的恋爱退没退场。}
\qarow{一成多不谈恋爱，剩下的多数没说不谈？}{对。还有超过四分之一说不确定；我方用它说明背景，主证据是嗑的人自己的话。}
\qarow{这份调查问的是嗑CP的人吗？}{不是，问的是未婚城市青年；嗑CP的人自己的话，来自许冠文的研究。}
\end{qa}
\end{stage}

\begin{stage}{7. 反三质询正二问题链}{90 秒，双边计时}
\textbf{总目标}：把“标准”变成“不用兑现”，把“宁缺毋滥”变成“缺”。
\begin{chain}{链一：兑现链}{让对方承认想象里的标准不需要兑现}
\q{在屏幕前要求专一，需要自己付出什么吗？}{答不需要：下一问；答需要：付出了什么？}
\q{不用付出的标准，能证明你相信它会来吗？}{答能：请评委记住；答不能：收束}
\close{感谢，对方承认屏幕前的标准不用兑现。}
\end{chain}
\begin{chain}{链二：宁缺链}{把宁缺毋滥落到现实}
\q{宁缺毋滥，落到现实是缺，还是有？}{答缺：下一问；答有：那为什么还单身？}
\q{一直缺着的人，是更信爱情，还是更怕爱情？}{答更信：请对方说出证据；答更怕：收束}
\q{中青报那项调查，样本多少，您知道吗？}{答不知道：请评委记住这个数字没有样本}
\close{感谢，宁缺毋滥落到现实，就是缺。}
\end{chain}
\end{stage}

\begin{stage}{8. 反一对辩要点}{双方各 90 秒}
\textbf{任务}：收拢前两段：把全场带回“更深”，谁是谁的来源。
\begin{points}{进攻点（4–5 个，每个 15–20 秒）}
\item 更深比的是来源：是因为不敢，才去嗑。对方说是因为想要，那想要为什么只敢放在屏幕那头？
\item 我方的问题还在：真信，为什么要放在一段碰不到、结局写好的关系里？
\item 许冠文那句话前半是很高的价值，后半是每一步都充满风险。决定人们怎么做的，是后半。
\item 画好终点才敢投入：不信结局会自己好。
\end{points}
\begin{points}{备用点（6–8 个）}
\item 对方说门口：一直站在门口不进去，是不信门里的人会等你。
\item 对方说补充：补充的前提是有正餐。
\item 对方说“又”相信：要靠别人的甜才能信一次，信在别人身上。
\item 对方说圈子窄：嗑CP的人把时间给了屏幕。
\item 对方说怕说明当真：当真又不敢，就是不信它靠得住。
\item 对方说宁缺毋滥：缺，是结果。
\end{points}
\begin{points}{对方最可能问的三个问题 → 回应}
\item 你方承认热潮里有信吗？ → 承认，信在别人身上；我方比的是谁在底下。
\item 不信的人为什么嗑好结局？ → 因为只有写好的结局才安全。
\item 旁观不等于怀疑吧？ → 不是所有旁观，是说出“现实太险”的旁观。
\end{points}
\end{stage}

\begin{kit}{9. 反三申论}{套件 · 组装后 380–403 字 / 90 秒}
\assembly{开头 → 回应正三申论（任选 1） → 处理正方最强点（任选 1） → 推进论点二 → 收尾}
\begin{fixed}{开头}
谢谢主席，各位。这一篇我做三件事：先回应正三，再处理正方最有力的一点，最后推进我方第二点。
\end{fixed}
\begin{pick}{回应正三申论}{1}
\begin{module}{如果正三申论主打“婚、育、恋里，恋爱被放下得最少”}
正三刚才说，婚育恋三样里，爱情被放下得最少。我方同意，可最少被放下，不等于被相信。意愿排第一，说明人们还想要；想要和相信之间，差的正是敢不敢。
\end{module}
\begin{module}{如果正三申论主打“对多数人是补充，不是替代”}
正三刚才说，嗑CP对多数人是补充。可补充的前提，是正餐还在吃。人们自己的恋爱，一成多说将来不谈，超过四分之一说不确定。这份补充，正在变成主食。
\end{module}
\begin{fallback}{正三申论主打不在以上两条时}
正三这一篇，我方先记下来，放回判准里看：它讲的是人们想要什么，还是人们为什么只敢这样要？判准要的是后一个，而这一个，请正方正面回答。
\end{fallback}
\end{pick}
\begin{pick}{处理正方最强点}{1}
\begin{module}{如果正方全场最有力的是“嗑的是专一、双向，不信的人不会嗑这些”}
正方全场最有力的，是这一条：不信爱情的人，为什么偏偏嗑专一和双向？我方的回答是：正因为不信它在现实里靠得住，才要在一个安全的地方看它成立。专一放在一段不会伤到自己的关系里，才敢要。想象里的标准越高，越说明现实里不敢试。
\end{module}
\begin{module}{如果正方全场最有力的是“怕风险，说明把爱当真”}
正方全场最有力的，是说怕风险说明当真。我方完全同意，人们把爱当真。可当真，和相信它靠得住，是两件事。当真，所以知道它会伤人；不信它靠得住，所以只敢在结局写好的地方爱。当真又不敢，这恰恰是中立定义里的怀疑。
\end{module}
\begin{fallback}{正方最有力的一点不在以上两条时}
不管正方最有力的是哪一条，都请评委用判准的两半去量：嗑什么，和为什么这样嗑。人们嗑专一、嗑双向，这些我方从一开始就不否认。内容说明人们想要什么，方式说明人们敢不敢。一个想要又不敢的人，底下是什么，评委心里有数。为什么偏偏旁观，这一问请正方回答。
\end{fallback}
\end{pick}
\begin{fixed}{推进论点二}
我方第二点，要的是保险，今天再补一层。法国哲学家巴迪欧在《爱的多重奏》里讲过一张婚恋网站的海报，上面写着：没有痛苦的完美爱情。他管这个叫爱面对的“安全威胁”，先把风险算尽，再去爱。嗑CP走的是同一条路：结局先写好，感情再放进去。共青团中央课题组二〇二一年的调查里，不想恋爱的青年，七成半选了“一个人很好，谈恋爱很麻烦”。请评委记住，麻烦两个字的意思是：付出了，未必换得回来。
\end{fixed}
\begin{fixed}{收尾}
说到底，正方要回答的只有一个问题：如果真信爱情，为什么非要先给它上一份保险？
\end{fixed}
\end{kit}

\begin{stage}{10. 反三被质询防守预案}{正一质询，90 秒双边计时}
\textbf{原则}：正一这一轮是为结辩取材，重点守住“承认有信，但信在上面”。
\begin{qa}
\qarow{嗑CP的人最想看到的结局，是好还是坏？}{好的。只有写好的好结局，他们才敢投入。}
\qarow{想看好结局，是相信它存在，还是不信？}{相信它在故事里存在；不信它在自己的现实里靠得住。}
\qarow{您方承认热潮里也有信吗？}{承认，信在上面，在屏幕那头。}
\qarow{是因为不信才嗑，还是因为想要才嗑？}{想要，但不敢，所以嗑。想要在上面，不敢在底下。}
\qarow{七成半嫌麻烦，麻烦是成本，不是怀疑吧？}{嫌麻烦，是觉得付出换不来值得的结果，这就是不信它值得投入。}
\qarow{巴迪欧批的是婚恋网站，不是嗑CP吧？}{对，他批的是一种逻辑：先保结局再去爱。嗑CP是同一种逻辑。}
\end{qa}
\end{stage}

\begin{stage}{11. 反一质询正三问题链}{90 秒，双边计时}
\textbf{总目标}：为反一结辩取材：对方承认人们不敢；对方把“不敢”和“不信”分开，分不清。
\begin{chain}{链一：不敢链}{让对方承认现实里人们不敢}
\q{现实中的亲密关系每一步都充满风险，您方同意吗？}{答同意：下一问；答不同意：许冠文原话，请对方反驳}
\q{因为风险大才去嗑，算不算不敢？}{答算：收束；答不算：那算什么？}
\close{感谢，对方承认人们因为风险而不敢。}
\end{chain}
\begin{chain}{链二：分界链}{逼对方给出“不敢”和“不信”的分界}
\q{不敢把爱放在自己身上，是不是不信它靠得住？}{答是：收束；答不是：下一问}
\q{不敢和不信，您方能说出区别吗？}{对方给出区别：记下，结辩里检验；答不出：请评委记住}
\cut{好，我换个问法：敢不敢，和靠不靠得住，有没有关系？}
\close{感谢，不敢和不信，对方没有分开。}
\end{chain}
\end{stage}

\begin{stage}{12. 反二对辩要点}{双方各 90 秒}
\textbf{任务}：结辩前最后一次交锋，守住判准的另一半：为什么用旁观的方式。不开新战场。
\begin{points}{进攻点（4–5 个，每个 15–20 秒）}
\item 判准有两半：嗑什么，为什么这样嗑。正方只讲了嗑什么。
\item 屏幕那头的爱，随时可以抽身、不用付出。只敢这样爱，是信还是疑？
\item 结局写好才敢投入，说明不信结局会自己好。
\item 许冠文那句话的后半：每一步都充满风险。决定人们怎么做的，是这一半。
\end{points}
\begin{points}{备用点（6–8 个）}
\item 对方说标准高：不用兑现的标准，谁都敢拉满。
\item 对方说门口：站在门口不进去，是不信门里值得。
\item 对方说怕说明当真：当真又不敢，就是不信它靠得住。
\item 对方说补充：一成多说不谈，超过四分之一说不确定，补充在变成主食。
\item 对方说宁缺毋滥：落到现实就是缺。
\item 对方说“又”相信：靠别人才能信，信在别人身上。
\end{points}
\begin{points}{对方最可能问的三个问题 → 回应}
\item 你方承认热潮里有信？ → 承认，信在上面；底下是不敢。
\item 麻烦不是怀疑吧？ → 嫌麻烦，是不信付出换得来值得的结果。
\item 看球的人不上场也不怀疑足球？ → 看球的人没说过球场太危险。
\end{points}
\end{stage}

\begin{kit}{13. 反一结辩}{套件 · 组装后 385–414 字 / 90 秒}
\assembly{归纳分歧（任选 1） → 处理最强点（任选 1） → 封路（任选 1） → 论点收束 → 价值升华 → 结尾}
\begin{pick}{归纳分歧}{1}
\begin{module}{如果全场主战场落在“嗑的内容说明什么”}
谢谢主席。今天吵得最久的，是嗑的内容。我方承认，内容说明人们想要什么。可这份想要，只敢放在屏幕那头，放在一段自己进不去的关系里。
\end{module}
\begin{module}{如果全场主战场落在“不敢算不算不信”}
谢谢主席。今天吵得最久的，是不敢算不算不信。不敢把爱放在自己身上，就是不信它靠得住。中立定义写得清楚：怀疑，包括不信它靠得住。
\end{module}
\begin{fallback}{主战场不在以上两条时}
谢谢主席。今天的分歧，其实就一句：热潮底下，是想要，还是不敢。我方说，想要在上面，不敢在底下。先有不敢，人们才去屏幕那头找想要的东西。
\end{fallback}
\end{pick}
\begin{pick}{处理最强点}{1}
\begin{module}{如果正方全场最有力的是“嗑的是专一和双向”}
正方最有力的，是人们嗑专一和双向。可专一放在一段不会伤到自己的关系里，才敢要。想象里的标准越高，越说明现实里不敢试。标准，是在屏幕那头升起来的；落到自己身上，它很少被拿出来用。
\end{module}
\begin{module}{如果正方全场最有力的是“疑的只是遇见”}
正方最有力的，是疑的只是遇见。可一份爱，扛不住圈子窄，扛不住热搜上的负面新闻，它在人们心里就是靠不住的。门口的疑站得久了，就会变成门里的疑：人们开始相信，门里那个人，大概不会为自己等。
\end{module}
\begin{fallback}{正方最有力的一点不在以上两条时}
正方最有力的那一点，请评委用判准的另一半去量：人们为什么只敢这样要？为什么偏偏放在屏幕那头？想要什么，我方不否认；敢不敢要，才是题目问的“更深”。
\end{fallback}
\end{pick}
\begin{pick}{封路}{1}
\begin{module}{如果正方全场主打“内容是理想，所以是信”}
正方待会儿还有最后一发。如果它仍然落在嗑的内容上，请评委替我们问一句：这份理想，有没有人敢放在自己身上？如果没有人敢，这份理想就还在屏幕那头，那它折射出的就不是信。
\end{module}
\begin{module}{如果正方全场主打“疑的是门口，门里有人”}
正方全场说，疑的是门口。如果最后一发还是门口，请评委替我们问一句：一直不进门的人，是嫌门口太挤，还是不信门里有人在等？站了这么久，答案其实已经在脚下。
\end{module}
\begin{fallback}{正方全场路线不在以上两条时}
正方还有最后一发。请评委拿着判准去听：它有没有回答，人们为什么只敢旁观，只敢在结局写好的地方爱？没有回答，那不管这一发多动人，热潮底下就还是疑。
\end{fallback}
\end{pick}
\begin{fixed}{论点收束}
我方两点，到最后都还站着。第一点，信在远处，爱只放在屏幕那头。第二点，要的是保险，结局写好了才敢投入。两点合起来就是一句话：信在上面，疑在底下，底下才是更深的那一层。
\end{fixed}
\begin{fixed}{价值升华}
你想想这样一个画面。凌晨一点，一个女生在弹幕里打下“他们一定要幸福啊”，然后关掉了屏幕。手机里还有一条消息，是一个对她有好感的人发来的，她一直没有回。不是不想，是怕。我们今天争的，不是年轻人爱不爱看甜。是那条没有回的消息。她相信屏幕里的人会幸福，却不相信自己也可以。
\end{fixed}
\begin{fixed}{结尾}
说到底，我们希望有一天，她敢回。谢谢。
\end{fixed}
\end{kit}
\begin{contingency}
\ifthen{时间不够}{先把处理最强点压成一句；价值升华与结尾不删。}
\end{contingency}
\end{document}
